\corrigesChapitre{11}{Énergie cinétique et bilan mécanique}
\par\noindent\begin{minipage}{\linewidth}
\begin{corrige}{1}{La forme d’énergie}
\begin{questions}
\item L’énergie cinétique.
\item En joules, J.
\end{questions}
\end{corrige}
\end{minipage}
\par\noindent\begin{minipage}{\linewidth}
\begin{corrige}{2}{Les bonnes unités}
\begin{questions}
\item kg et m/s.
\item $0{,}500\unite{kg}$.
\end{questions}
\end{corrige}
\end{minipage}
\par\noindent\begin{minipage}{\linewidth}
\begin{corrige}{3}{La vitesse}
\begin{questions}
\item $10\unite{m/s}$.
\item $20\unite{m/s}$.
\end{questions}
\end{corrige}
\end{minipage}
\par\noindent\begin{minipage}{\linewidth}
\begin{corrige}{4}{Un calcul direct}
\begin{questions}
\item $\tfrac12\times2\times9=9\unite J$.
\item $0\unite J$.
\end{questions}
\end{corrige}
\end{minipage}
\par\noindent\begin{minipage}{\linewidth}
\begin{corrige}{5}{Deux facteurs}
\begin{questions}
\item Ec double.
\item Ec est multipliée par quatre.
\end{questions}
\pcFigureEx{figures/rev-cor-ec-facteurs.pdf}{Courbe de rapports à masse constante : le facteur appliqué à la vitesse est mis au carré pour l’énergie cinétique.}
\end{corrige}
\end{minipage}
\par\noindent\begin{minipage}{\linewidth}
\begin{corrige}{6}{Le ballon}
\begin{questions}
\item $E_c=0{,}5\times0{,}400\times25=5\unite J$.
\item $20\unite J$.
\end{questions}
\pcFigureEx{figures/rev-cor-ec-facteurs.pdf}{Courbe de rapports à masse constante : le facteur appliqué à la vitesse est mis au carré pour l’énergie cinétique.}
\end{corrige}
\end{minipage}
\par\noindent\begin{minipage}{\linewidth}
\begin{corrige}{7}{Une voiture}
\begin{questions}
\item $v=15\unite{m/s}$ ; $E_c=112\,500\unite J$.
\item $112{,}5\unite{kJ}$.
\end{questions}
\par\smallskip\noindent\textbf{Pour comprendre.}\enspace $54\unite{km/h}=54/3{,}6=15\unite{m/s}$. Alors $E_c=\frac12\times1\,000\times15^2=112\,500\unite J=112{,}5\unite{kJ}$. Mettre 54 au carré sans conversion donnerait une valeur fausse.\par
\end{corrige}
\end{minipage}
\par\noindent\begin{minipage}{\linewidth}
\begin{corrige}{8}{Ralentir}
\begin{questions}
\item $3\,000\unite J$ puis $750\unite J$.
\item $2\,250\unite J$, transférés ou convertis, notamment en thermique.
\end{questions}
\par\smallskip\noindent\textbf{Pour comprendre.}\enspace $E_{c,i}=\frac12\times60\times10^2=3\,000\unite J$, puis $E_{c,f}=\frac12\times60\times5^2=750\unite J$. L’écart est $3\,000-750=2\,250\unite J$.\par
\pcFigureEx{figures/rev-cor-freinage.pdf}{L’écart de 2 250 J ne disparaît pas : il est transféré ou converti, notamment par les frottements.}
\end{corrige}
\end{minipage}
\par\noindent\begin{minipage}{\linewidth}
\begin{corrige}{9}{Une chute idéale}
\begin{questions}
\item $E_c=40-15=25\unite J$.
\item Les frottements réels peuvent transférer une partie de l’énergie mécanique vers le milieu.
\end{questions}
\pcFigureEx{figures/rev-cor-chute.pdf}{Dans le modèle sans frottements, l’énergie mécanique reste 40 J ; les deux formes se redistribuent.}
\end{corrige}
\end{minipage}
\par\noindent\begin{minipage}{\linewidth}
\begin{corrige}{10}{Le rapport}
\begin{questions}
\item Rapport $3^2=9$.
\item Rapport $120/80=1{,}5$.
\end{questions}
\pcFigureEx{figures/rev-cor-ec-facteurs.pdf}{Courbe de rapports à masse constante : le facteur appliqué à la vitesse est mis au carré pour l’énergie cinétique.}
\end{corrige}
\end{minipage}
\par\noindent\begin{minipage}{\linewidth}
\begin{corrige}{11}{Une hausse de 20 \%}
\begin{questions}
\item $1{,}2^2=1{,}44$.
\item 44 \%.
\end{questions}
\par\smallskip\noindent\textbf{Pour comprendre.}\enspace Une hausse de 20 \% donne le facteur de vitesse $1+20/100=1{,}2$. Le facteur d’énergie est $1{,}2^2=1{,}44$ : la hausse d’énergie vaut donc $(1{,}44-1)\times100=44\,\%$.\par
\pcFigureEx{figures/rev-cor-ec-facteurs.pdf}{Courbe de rapports à masse constante : le facteur appliqué à la vitesse est mis au carré pour l’énergie cinétique.}
\end{corrige}
\end{minipage}
\par\noindent\begin{minipage}{\linewidth}
\begin{corrige}{12}{Changer de référentiel}
\begin{questions}
\item Oui, puisqu’il se déplace par rapport au sol.
\item Le référentiel et la vitesse du système dans ce référentiel.
\end{questions}
\end{corrige}
\end{minipage}
\par\noindent\begin{minipage}{\linewidth}
\begin{corrige}{13}{La force n’est pas l’énergie}
\begin{questions}
\item 250 J si le calcul est correct.
\item Non ; il faut un modèle et des données sur le freinage, la route et les autres actions.
\end{questions}
\end{corrige}
\end{minipage}
\par\noindent\begin{minipage}{\linewidth}
\begin{corrige}{14}{Un bilan réel}
\begin{questions}
\item Vers de l’énergie thermique du système et du milieu, notamment par frottement.
\item Non, à condition d’élargir le bilan aux transferts et aux autres formes d’énergie.
\end{questions}
\pcFigureEx{figures/rev-cor-chute-reelle.pdf}{Le bilan mécanique diminue de 15 J ; le bilan total inclut l’énergie transférée au milieu.}
\end{corrige}
\end{minipage}
\par\noindent\begin{minipage}{\linewidth}
\begin{corrige}{15}{Comparer deux véhicules}
\begin{questions}
\item A : 160 000 J ; B : 80 000 J. A a deux fois plus d’énergie.
\item Le carré de la vitesse compense ici plus que la différence de masse.
\end{questions}
\par\smallskip\noindent\textbf{Pour comprendre.}\enspace A : $\frac12\times800\times20^2=160\,000\unite J$. B : $\frac12\times1\,600\times10^2=80\,000\unite J$. Le rapport A/B vaut 2 : la masse seule ne suffit pas à comparer.\par
\pcFigureEx{figures/rev-cor-vehicules.pdf}{Comparer des résultats dans le même référentiel, après avoir pris le carré des vitesses.}
\end{corrige}
\end{minipage}
\par\noindent\begin{minipage}{\linewidth}
\phantomsection\label{cor-defi-11}\pcSolutionDefi{11}{Une étude de sécurité}\begin{questions}
\item À 50 km/h : $v\approx13{,}89\unite{m/s}$ et $E_c\approx1{,}16\times10^5\unite J$. À 100 km/h : environ $4{,}63\times10^5\unite J$.
\item Même masse et vitesse doublée : énergie quadruplée.
\item Une vitesse plus grande augmente fortement l’énergie à dissiper. Le calcul suppose le même ensemble et le même référentiel. Il ne donne pas directement la distance d’arrêt. Celle-ci dépend aussi du temps de réaction, de l’adhérence et du freinage.
\end{questions}
\par\smallskip\noindent\textbf{Pour comprendre.}\enspace Conserver les valeurs non arrondies de $50/3{,}6$ et $100/3{,}6$ jusqu’au calcul final. Le rapport des vitesses est exactement 2, donc celui des énergies exactement 4 dans le modèle, même si les résultats affichés sont arrondis.\par
\pcFigureEx{figures/rev-cor-ec-facteurs.pdf}{Courbe de rapports à masse constante : le facteur appliqué à la vitesse est mis au carré pour l’énergie cinétique.}
\end{minipage}
\par\noindent\begin{minipage}{\linewidth}
\phantomsection\label{cor-activite-11}\pcSousTitre{Pistes pour l’investigation}
À masse fixée, Ec est proportionnelle à v² : le graphe Ec(v) est courbe et Ec(v²) est une droite passant par l’origine dans le modèle. En chute idéale, la baisse d’énergie potentielle est compensée par la hausse d’énergie cinétique. Avec frottements, ajouter les transferts thermiques et préciser les limites du système pour conserver l’énergie totale.\par\end{minipage}
